\chapter{Social Bots als Medien}
\label{ch:01_social-bots-as-media}

\begin{universityTask}
    Offensichtlich sind Social Bots -- ähnlich wie Chatbots -- interaktive Systeme und damit quarternäre Medien. Allerdings werden sie in der Regel nicht direkt, sondern wiederum innerhalb eines weiteren Mediums eingesetzt, durch das nicht zuletzt ihre eigene Medialität sozusagen versteckt wird. Der Einfachheit halber wollen wir hier annehmen, dass es sich dabei um ein Nachrichtenforum oder um einen Chatroom handelt.

    \vspace{0.5cm}

    \begin{enumerate}[label=(\alph*)]
        \item Geben Sie zunächst an, welche Kommunikationspartner bei einer solchen Social-Bot-Nachricht überhaupt sinnvollerweise medienwissenschaftlich zu betrachten sind.
            \universityPoints{2}
        \item Was zeichnet die quarternärmediale Kommunikation gegenüber sekundär- und tertiärmedialen Kommunikationen aus? Bitte geben Sie dazu die Kriterien für diese drei Medienklassen an den Beispielen Social Bot, Chatroom und Online-Forum an.
            \universityPoints{6}
        \item Vergleichen Sie die Bot-Nachrichten (i) mit den \q{normalen} Nachrichten in einem Chatroom (ii) und den Artikeln in einer Tageszeitung (iii) hinsichtlich (a) Authentizität, (b) Integrität, (c) Vertraulichkeit und (d) Nicht-Abstreitbarkeit. Welche Eigenschaft ist in welchem Medium intrinsisch verfügbar (Tabelle mit Begründungen!)?
            \universityPoints{12}
        \item Wie und wieso lässt sich der Begriff \q{Avatar} auf die Verwendung von Social Bots beziehen?
            \universityPoints{2}
    \end{enumerate}
\end{universityTask}

Dieser Text enthält die Lösung der obigen Aufgabenstellung.